%% Copyleft 2018 Jean-Baptiste Louvet
%% Copyright (C) 2014 Dorian Depriester
%% http://blog.dorian-depriester.fr
%%
%% This file may be distributed and/or modified under the conditions
%% of the LaTeX Project Public License, either version 1.3c of this
%% license or (at your option) any later version. The latest version
%% of this license is in:
%%
%%    http://www.latex-project.org/lppl.txt
%%
%% and version 1.3c or later is part of all distributions of LaTeX
%% version 2006/05/20 or later.
%%
%% This work has the LPPL maintenance status `maintained'.
%%
%% The Current Maintainer of this work is Dorian Depriester
%% <contact [at] dorian [-] depriester [dot] fr>.
%%
%% This is main.tex for French PhD Thesis.
%% See http://blog.dorian-depriester.fr/latex/template-these/template-complet-pour-manuscrit-de-these for help


%%%%%%%%%%%%%%%%%%%%%%%%%%%%%%%%%%%%%%%%%
%           Fichier maitre				%
%%%%%%%%%%%%%%%%%%%%%%%%%%%%%%%%%%%%%%%%%

% \documentclass[a4paper, french, 11pt, final, twoside,openright]{report}
\documentclass[a4paper, 12pt, twoside,openright]{report}
\usepackage[french,english]{babel}
\usepackage[onehalfspacing]{setspace}
\usepackage{amssymb}
\usepackage{algorithm}
\usepackage{amsmath}
\usepackage{amsfonts}
\usepackage{float}
\usepackage{diagbox}
\usepackage{algpseudocode}
\usepackage[cal=pxtx]{mathalfa}
\usepackage[dvipsnames]{xcolor}
\usepackage{psfrag,overpic}
%\usepackage{comment}
\renewcommand{\algorithmicrequire}{\textbf{Input:}}
\renewcommand{\algorithmicensure}{\textbf{Output:}}

%\usepackage{biblatex}
%\addbibresource{Manuscript.bib}

\renewcommand{\labelenumi}{\arabic{enumi})}

%% Réglage des fontes et typo    
\usepackage[utf8]{inputenc}		% LaTeX, comprend les accents !
\usepackage[T1]{fontenc}

% \usepackage[square,sort&compress,sectionbib]{natbib}		% Doit être chargé avant babel
% \usepackage{chapterbib}
% 	\renewcommand{\bibsection}{\section{Références}}		% Met les références biblio dans un \section (au lieu de \section*)

\usepackage{amsthm}		
\usepackage{babel}
\usepackage{lmodern}
\usepackage{ae,aecompl}										% Utilisation des fontes vectorielles modernes
\usepackage[upright]{fourier}
\usepackage{xcolor} % Coloration de texte
\usepackage{csquotes}



%%%%%%%%%%%%%%%%%%%%%%%%%%%%%%%%%%%%%%%%%%%%%%%%%%%%%%%%%%%%%%%%%%%%%
% Allure générale du document
\usepackage{enumerate}
\usepackage{enumitem}
\usepackage[section]{placeins}	% Place un FloatBarrier à chaque nouvelle section
\usepackage{epigraph}
\usepackage[font={small}]{caption}
\usepackage[english,nohints]{minitoc}		% Mini table des matières, en français
	\setcounter{minitocdepth}{2}	% Mini-toc détaillées (sections/sous-sections)
\usepackage[notbib]{tocbibind}		% Ajoute les Tables	des Matières/Figures/Tableaux à la table des matières

\usepackage{datetime} % Pour avoir l'heure de compilation
\usepackage{xspace}
\usepackage{fancyhdr}

%%%%%%%%%%%%%%%%%%%%%%%%%%%%%%%%%%%%%%%%%%%%%%%%%%%%%%%%%%%%%%%%%%%%%
%% Maths                         
\usepackage{amsmath}			% Permet de taper des formules mathématiques
\usepackage{amssymb}			% Permet d'utiliser des symboles mathématiques
\usepackage{amsfonts}			% Permet d'utiliser des polices mathématiques
\usepackage{nicefrac}			% Fractions 'inline'


%%%%%%%%%%%%%%%%%%%%%%%%%%%%%%%%%%%%%%%%%%%%%%%%%%%%%%%%%%%%%%%%%%%%%
%% Tableaux
\usepackage{multirow}
\usepackage{booktabs}
\usepackage{colortbl}
\usepackage{tabularx}
\usepackage{multirow}
\usepackage{threeparttable}
\usepackage{etoolbox}
	\appto\TPTnoteSettings{\footnotesize}
\addto\captionsfrench{\def\tablename{{\textsc{Tableau}}}}	% Renome 'table' en 'tableau'
\usepackage{colortbl}

%%%%%%%%%%%%%%%%%%%%%%%%%%%%%%%%%%%%%%%%%%%%%%%%%%%%%%%%%%%%%%%%%%%%%
%% Eléments graphiques                    
\usepackage{graphicx}			% Permet l'inclusion d'images
\usepackage{subcaption}
\usepackage{pdfpages}
\usepackage{rotating}
\usepackage{pgfplots}
	\usepgfplotslibrary{groupplots}
\usepackage{tikz}
	\usetikzlibrary{backgrounds,automata}
	\pgfplotsset{width=7cm,compat=1.3}
	\tikzset{every picture/.style={execute at begin picture={
   		\shorthandoff{:;!?};}
	}}
	\pgfplotsset{every linear axis/.append style={
		/pgf/number format/.cd,
		use comma,
		1000 sep={\,},
	}}
\usepackage{eso-pic}
\usepackage{import}

%%%%%%%%%%%%%%%%%%%%%%%%%%%%%%%%%%%%%%%%%%%%%%%%%%%%%%%%%%%%%%%%%%%%%
%% Mise en forme du texte        
\usepackage{xspace}
\usepackage[load-configurations = abbreviations]{siunitx}
	\DeclareSIUnit{\MPa}{\mega\pascal}
	\DeclareSIUnit{\micron}{\micro\meter}
	\DeclareSIUnit{\tr}{tr}
	\DeclareSIPostPower\totheM{m}
	\sisetup{
	locale = FR,
	  inter-unit-separator=$\cdot$,
	  range-phrase=~\`{a}~,     	% Utilise le tiret court pour dire "de... à"
	  range-units=single,  		% Cache l'unité sur la première borne
	  }

\usepackage[version=3]{mhchem}	% Equations chimiques
\usepackage{textcomp}
\usepackage{array}
\usepackage{hyphenat}

%%%%%%%%%%%%%%%%%%%%%%%%%%%%%%%%%%%%%%%%%%%%%%%%%%%%%%%%%%%%%%%%%%%%%
%% Navigation dans le document
\usepackage[pdftex,hidelinks,
			%pagebackref=true,
			]{hyperref}	% Créera automatiquement les liens internes au PDF
					% Doit être chargé en dernier (Sauf exceptions ci-dessous)
			

%%%%%%%%%%%%%%%%%%%%%%%%%%%%%%%%%%%%%%%%%%%%%%%%%%%%%%%%%%%%%%%%%%%%%
%% Packages qui doivent être chargés APRES hyperref	             
\usepackage[top=2.5cm, bottom=2.5cm, left=3.5cm, right=3.cm, includeheadfoot = true, 
			headheight=15pt]{geometry}

\usepackage{fancyhdr}			% Entête et pieds de page. Doit être placé APRES geometry
	\pagestyle{fancy}		% Indique que le style de la page sera justement fancy
	\lfoot[\thepage]{} 		% gauche du pied de page
	\cfoot{} 			% milieu du pied de page
	\rfoot[]{\thepage} 		% droite du pied de page
	\fancyhead[RO, LE] {}	
	


%%%%%%%%%%%%%%%%%%%%%%%%%%%%%%%%%%%%%%%%%%%%%%%%%%%%%%
%% Pour aider à l'écriture de la thèse et au débugage
%%%%%%%%%%%%%%%%%%%%%%%%%%%%%%%%%%%%%%%%%%%%%%%%%%%%%%
\usepackage[french,color=white,linecolor=black]{todonotes} % Voir http://mirror.ibcp.fr/pub/CTAN/macros/latex/contrib/todonotes/todonotes.pdf
%\usepackage{showkeys} % Voir http://mirrors.ircam.fr/pub/CTAN/macros/latex/required/tools/showkeys.pdf






%%%%%%%%%%%%%%%%%%%%%%%%%%%%%%%%%%%%%%%%%%%%%%%%%%%%%%%%%%%%%%%%%%%%%
%% Tableaux

% my own settings
\definecolor{cmt}{rgb}{1,.5,0} % color for comments and marks
\definecolor{red}{rgb}{1,0,0}
\definecolor{green}{rgb}{0,.6,0}

\usepackage{algorithm,algorithmicx,algpseudocode}
\usepackage{booktabs} % To thicken table lines
\usepackage{amsmath} % For math equation management
\usepackage{pifont} % For check and cross marks
\usepackage{ulem} % For underlines that can break in line
\usepackage[toc,page]{appendix} % appendix

\newcommand{\X}{\ensuremath{\mathcal{X}}} % kernel space
\newcommand{\Hil}{\ensuremath{\mathcal{H}}} % Hilbert space
\newcommand{\bigO}{\ensuremath{\mathcal{O}}} % 
\newcommand{\R}{\ensuremath{\mathbb{R}}} % real domain
\newcommand{\yes}{\textcolor{green}{\textbf{\ding{51}}}} % check marks
\newcommand{\xmark}{\textcolor{red}{\textbf{\ding{55}}}} % cross marks
\newcommand{\mmark}[1]{\textcolor{cmt}{\textit{\textbf{\uline{#1}}}}} % format for marks
\newcommand{\mr}[2]{\multirow{#1}*{#2}} % multirow
\newcommand{\mc}[2]{\multicolumn{#1}{c}{#2}} % multicolumn
	
\newtheorem{definition}{Definition} 

\algnewcommand{\algorithmicforeach}{\textbf{for each}}
\algdef{SE}[FOR]{ForEach}{EndForEach}[1]
  {\algorithmicforeach\ #1\ \algorithmicdo}% \ForEach{#1}
  {\algorithmicend\ \algorithmicforeach}% \EndForEach
 \let\oldReturn\Return
\renewcommand{\Return}{\State\oldReturn}

%%%%%%%%%%%%%%%%%%%%%%%%%%%%%%%%%%%%%%%%%%%%%%%%%%%%%%
%% Pour aider à l'écriture de la thèse et au débugage
%%%%%%%%%%%%%%%%%%%%%%%%%%%%%%%%%%%%%%%%%%%%%%%%%%%%%%
\usepackage[english,color=white,linecolor=black]{todonotes} % Voir http://mirror.ibcp.fr/pub/CTAN/macros/latex/contrib/todonotes/todonotes.pdf
%\usepackage{showkeys} % Voir http://mirrors.ircam.fr/pub/CTAN/macros/latex/required/tools/showkeys.pdf
%%%%%%%%%%%%%%%%%%%%%%%%%%%%
%% Packages perso
%%%%%%%%%%%%%%%%%%%%%%%%%%%%
		% Liste des packages et de leurs options
%%%%%%%%%%%%%%%%%%%%%%%%%%%%%%%%%%%%%%%%
%           Commandes perso            %
%%%%%%%%%%%%%%%%%%%%%%%%%%%%%%%%%%%%%%%%

%% Figures centrées, et en position 'here, top, bottom or page'
\newenvironment{figureth}{%
		\begin{figure}[htbp]
			\centering
	}{
		\end{figure}
		}
		
		
%% Tableaux centrés, et en position 'here, top, bottom or page'
\newenvironment{tableth}{%
		\begin{table}[htbp]
			\centering
			%\rowcolors{1}{coleurtableau}{coleurtableau}
	}{
		\end{table}
		}

%% Sous-figures centrées, en position 'top'		
\newenvironment{subfigureth}[1]{%
	\begin{subfigure}[t]{#1}
	\centering
}{
	\end{subfigure}
}

\newcommand{\citationChap}[2]{%
	\epigraph{\og \textit{#1} \fg{}}{#2}
}

%% On commence par une page impaire quand on change le style de numérotation de pages 
\let\oldpagenumbering\pagenumbering
\renewcommand{\pagenumbering}[1]{%
	\cleardoublepage
	\oldpagenumbering{#1}
}
	% Commandes et environnements perso
%% Copyleft 2018 Jean-Baptiste Louvet
%% Copyright (C) 2014 Dorian Depriester
%% http://blog.dorian-depriester.fr
%%
%% This file may be distributed and/or modified under the conditions
%% of the LaTeX Project Public License, either version 1.3c of this
%% license or (at your option) any later version. The latest version
%% of this license is in:
%%
%%    http://www.latex-project.org/lppl.txt
%%
%% and version 1.3c or later is part of all distributions of LaTeX
%% version 2006/05/20 or later.
%%
%% This work has the LPPL maintenance status `maintained'.
%%
%% The Current Maintainer of this work is Dorian Depriester
%% <contact [at] dorian [-] depriester [dot] fr>.
%%
%% This is PageDeGarde.tex for French PhD Thesis.

%%%%%%%%%%%%%%%%%%%%%%%%%%%%%%%%%%%%%%%%%%
%           Page de Garde		         %
%%%%%%%%%%%%%%%%%%%%%%%%%%%%%%%%%%%%%%%%%%

\makeatletter
\def\@ecole{école}
\newcommand{\ecole}[1]{
\def\@ecole{#1}
}

\def\@specialite{Spécialité}
\newcommand{\specialite}[1]{
  \def\@specialite{#1}
}

\def\@directeur{directeur}
\newcommand{\directeur}[1]{
  \def\@directeur{#1}
}

\def\@encadrant{}
\newcommand{\encadrant}[1]{
  \def\@encadrant{et #1}
}

\def\@version{}
\newcommand{\version}[1]{
  \def\@version{V #1 – compilé le \today{} à \currenttime}
}

\def\@jurya{}{}{}
\newcommand{\jurya}[3]{
  \def\@jurya{#1	& #2	& #3\\\hline}
}
\def\@juryb{}{}{}
\newcommand{\juryb}[3]{
  \def\@juryb{#1	& #2	& #3\\\hline}
}
\def\@juryc{}{}{}
\newcommand{\juryc}[3]{
  \def\@juryc{#1	& #2	& #3\\\hline}
}
\def\@juryd{}{}{}
\newcommand{\juryd}[3]{
  \def\@juryd{#1	& #2	& #3\\\hline}
}
\def\@jurye{}{}{}
\newcommand{\jurye}[3]{
  \def\@jurye{#1	& #2	& #3\\\hline}
}
\def\@juryf{}{}{}
\newcommand{\juryf}[3]{
  \def\@juryf{#1	& #2	& #3\\\hline}
}
\def\@juryg{}{}{}
\newcommand{\juryg}[3]{
  \def\@juryg{#1	& #2	& #3\\\hline}
}
\def\@juryh{}{}{}
\newcommand{\juryh}[3]{
  \def\@juryh{#1	& #2	& #3\\\hline}
}
\def\@juryi{}{}{}
\newcommand{\juryi}[3]{
  \def\@juryi{#1	& #2	& #3\\\hline}
}
\newcommand{\pagedegarde}{
\newgeometry{top=1cm, bottom=1cm, left=1cm, right=1cm}
  \begin{titlepage}
  \centering
      \includegraphics[width=0.3\textwidth]{Normandie-Universite.eps}\\
    \vspace{0.5cm}
	\colorbox{violet!75!black}{\begin{minipage}{\textwidth}
		\begin{center}
			\textcolor{white}{
				{\bf\huge THÈSE}
			}
		\end{center}
	\end{minipage}}\\
    \vspace{0.7cm}
        {\bf\Large Pour obtenir le diplôme de doctorat}\\
    \vspace{0.5cm}
   	{\bf\large Spécialité \@specialite}\\
    \vspace{0.5cm}
    	{\bf\large Préparée au sein de \@ecole}\\
    \vspace{0.7cm}
	%\fboxsep4pt
	\colorbox{black!30}{\begin{minipage}{\textwidth}
		\begin{center}
			{\bf\LARGE \@title}
		\end{center}
	\end{minipage}}

%         \colorbox{black!30}{
%         \makebox[\textwidth]{
%         \centering
%         \begin{tabular*}{@{}l@{}}
%         
%         {\bf\LARGE \@title}
%         
%         \end{tabular*}
%         }
%         }
        
%         \colorbox{black!30}{\parbox[height=2cm]{\textwidth}{\center\bf\LARGE \@title}}\\
    \vspace{0.7cm}
    	\textit{\bf\Large Présentée et soutenue par}\\
    \vspace{0.25cm}
    	{\Large {\bfseries \@author}} \\
    \vspace{0.7cm}
	\colorbox{orange!60}{\begin{minipage}{\textwidth}
		\begin{center}
			{\bf\large Thèse soutenue publiquement le \@date\\
			devant le jury composé de}\\
		\end{center}
	\end{minipage}}
% 	\begin{tabular}{!{\color{orange!60}\vrule}l!{\color{orange!60}\vrule}l!{\color{orange!60}\vrule}r!{\color{orange!60}\vrule}}
%                 \arrayrulecolor{orange!60}
% 		\@jurya\hline
% 		\@juryb\hline
% 		\@juryc\hline
% 		\@juryd\hline
% 		\@jurye\hline
% % 		\@juryf\hline
% % 		\@juryg\hline
% % 		\@juryh\hline
% % 		\@juryi\hline
% 	\end{tabular}\\
%     \vspace{-0.1cm}
        {\renewcommand{\arraystretch}{1.25}%
	\begin{tabular}[c]{!{\color{orange!60}\vrule}m{0.3\textwidth}!{\color{orange!60}\vrule}m{0.3415\textwidth}!{\color{orange!60}\vrule}m{0.3\textwidth}!{\color{orange!60}\vrule}}
% 	\begin{tabular}{m{0.3\textwidth}m{0.3415\textwidth}m{0.3\textwidth}}
	\arrayrulecolor{orange!60}
	\hline
		\@jurya
        \@juryb
		\@juryc
		\@juryd
		\@jurye
		\@juryf
		\@juryg
		\@juryh
		\@juryi
	\end{tabular}\\
	}
    \vspace{0.7cm}
%	{\bf\large Thèse dirigée par \@directeur{} \@encadrant}\\
%    \vfill
%	\vspace{0.7cm}
\vspace{0.7cm}
\vspace{0.7cm}
\vspace{0.7cm}
	
	\includegraphics[height=3cm]{INSA_Short.pdf}
	\hfill
	\includegraphics[height=3cm]{MIIS.png}
	\hfill
	\includegraphics[height=3cm]{LITIS.pdf}
  \end{titlepage}




\restoregeometry  
  
  
}
\makeatother


%\DeclareMathOperator{\sign}{sign}

% Méta-données du PDF

%\hypersetup{
%     pdfauthor={Auteur},
%     pdfsubject={Manuscrit de thèse de doctorat},
%     pdftitle={La meilleure manière d'aller d'un point $a$ à un point $b$},
%     pdfkeywords={TODO}%TODO
% }


%%%%%%%%%%%%%%%%%%%%%%%%%%%%%%%%%%%%%%%%%%%%%%%%%%%%%%%%%%%%%%%%%
%%   			Liste des fichiers à compiler					%
%%%%%%%%%%%%%%%%%%%%%%%%%%%%%%%%%%%%%%%%%%%%%%%%%%%%%%%%%%%%%%%%%
%	\includeonly{Chapitre1,Chapitre2,Annexes}


% Répertoire contenant les images
\graphicspath{{Images/}}
% Infos de la page de garde
\author{Qiushi CAO}
%¨\title{ Algorithms for $\ell_0$ based sparse representation}
\title{Semantic Technologies for the Modeling of Predictive Maintenance for a SME Network in the Framework of Industry 4.0}
\specialite{Informatique}
\ecole{INSA ROUEN NORMANDIE}
\date{26/06/2020}
%\directeur{Cecilia ZANNI-MERK}
%\encadrant{Christoph REICH. Encadrement: Ahmed SAMET} %Optionnel
\jurya{M. Thomas \textsc{GUYET}}{MCF HDR, AGROCAMPUS-OUEST/IRISA}{Rapporteur}
\juryb{Mme Nada \textsc{MATTA}}{PR, Université Technologique de Troyes}{Rapporteur}
\juryc{Mme Anne \textsc{H{\AA}KANSSON}}{PR, UiT - The Arctic University of Norway}{Examinateur}
\juryd{Mme Samira \textsc{SI-SAID CHERFI}}{PR, CNAM}{Examinateur}
\jurye{Mme Cecilia \textsc{ZANNI-MERK}}{PR, INSA Rouen Normandie}{Directeur de th\`ese}
\juryf{M. Christoph \textsc{REICH}}{PR, Hoschshule Furtwangen}{Co-directeur de th\`ese}%Codirectrice de thèse
\juryg{M. Ahmed \textsc{SAMET}}{MCF, INSA Strasbourg}{Encadrant de thèse}
\juryh{M. Fran\c cois \textsc{DE BERTRAND DE BEUVRON}}{MCF, INSA Strasbourg}{Encadrant de thèse}

%\newcommand{\ver}{0.0}% Version du document
%\version{\ver}% À commenter si l'on ne veut pas afficher la version et date de compilation sur la page de garde

%%%%%%%%%%%%%%%%%%%%%%%%%%%%%%%%%%%%%%%%%%%%%%%%%
% Ajoute un timestamp de compilation dans le footer du document 
% À commenter dans la version finale
%\fancyfoot[L]{}
\pagestyle{fancy}

\usepackage{tabulary}
\newcolumntype{C}[1]{>{\centering\arraybackslash}p{#1}} 

\DeclareMathSymbol{\R}{\mathalpha}{AMSb}{"52}
\DeclareMathSymbol{\N}{\mathalpha}{AMSb}{"4E}
\DeclareMathSymbol{\Z}{\mathalpha}{AMSb}{"5A}
\DeclareMathOperator*{\argmin}{arg\,min}
\DeclareMathOperator*{\argmax}{arg\,max}

\newcommand{\x}{\mathbf{x}}
\newcommand{\y}{\mathbf{y}}
\newcommand{\dd}{\mathbf{d}}
\newcommand{\e}{\mathbf{e}}
\newcommand{\z}{\mathbf{z}}
\newcommand{\bv}{\mathbf{v}}
\newcommand{\br}{\mathbf{r}}
\newcommand{\bu}{\mathbf{u}}
\newcommand{\s}{\mathbf{s}}
\newcommand{\bc}{\mathbf{c}}
\newcommand{\bb}{\mathbf{b}}
\newcommand{\balpha}{\pmb{\alpha}}
\newcommand{\ba}{\mathbf{a}}

\algdef{SE}[DOWHILE]{Do}{doWhile}{\algorithmicdo}[1]{\algorithmicwhile\ #1}
\newtheorem{theorem}{Theorem}[section]
\newtheorem{prop}{Proposition}
\newtheorem{defn}{Definition}
\newtheorem{remark}{Remark}


\newcommand{\prox}{\text{prox}}
\newcommand{\tr}{\text{tr}}
\newcommand{\diag}{\text{diag}}
\newcommand{\eig}{\text{eig}}
\newcommand{\sign}{\text{sign}}

\sloppy

\begin{document}
% Préambule
	\pagenumbering{roman}
	{\footnotesize
	\pagedegarde}
	\cleardoublepage

		% Table des matières
			\setcounter{tocdepth}{2}	% Pas besoin de trop détailler le sommaire ici (chapitres/sections/sous-sections)
			\chapter*{Acknowledgements}






			\chapter*{}
\normalsize
\vspace{-6cm}
\begin{center}
\large \textbf{Résumé}
\end{center}
\small

$\quad  \!\!\!\!\!\!\!\!\!\!\!$ Dans le domaine de la fabrication, la détection d'anomalies telles que les défauts et les défaillances mécaniques permet de lancer des tâches de maintenance prédictive, qui visent à prévoir les défauts, les erreurs et les défaillances futurs et à permettre des actions de maintenance. Avec la tendance de l'industrie 4.0, les tâches de maintenance prédictive bénéficient de technologies avancées telles que les systèmes cyberphysiques (CPS), l'Internet des objets (IoT) et l'informatique dématérialisée (cloud computing). Ces technologies avancées permettent la collecte et le traitement de données de capteurs qui contiennent des mesures de signaux physiques de machines, tels que la température, la tension et les vibrations.

Cependant, en raison de la nature hétérogène des données industrielles, les connaissances extraites des données industrielles sont parfois présentées dans une structure complexe. Des méthodes formelles de représentation des connaissances sont donc nécessaires pour faciliter la compréhension et l'exploitation des connaissances. En outre, comme les CPSs sont de plus en plus axées sur la connaissance, une représentation uniforme de la connaissance des ressources physiques et des capacités de raisonnement pour les tâches analytiques est nécessaire pour automatiser les processus de prise de décision dans les CPSs. Ces problèmes constituent des obstacles pour les opérateurs de machines qui doivent effectuer des opérations de maintenance appropriées.

Pour relever les défis susmentionnés, nous proposons dans cette thèse une nouvelle approche sémantique pour faciliter les tâches de maintenance prédictive dans les processus de fabrication. En particulier, nous proposons quatre contributions principales: i) un cadre ontologique à trois niveaux qui est l'élément central d'un système de maintenance prédictive basé sur la connaissance; ii) une nouvelle approche sémantique hybride pour automatiser les tâches de prédiction des pannes de machines, qui est basée sur l'utilisation combinée de chroniques (un type plus descriptif de modèles séquentiels) et de technologies sémantiques; iii) a new approach that uses clustering methods with Semantic Web Rule Language (SWRL) rules to assess failures according to their criticality levels; iv) une nouvelle approche d'affinement de la base de règles qui utilise des mesures de qualité des règles comme références pour affiner une base de règles dans un système de maintenance prédictive basé sur la connaissance. Ces approches ont été validées sur des ensembles de données réelles et synthétiques.

Mots-clés : Industrie 4.0, Maintenance prédictive, Prévision des défaillances de machines, Évaluation de la criticité des défaillances, Ingénierie de l'ontologie, Rule-based Reasoning, Rule Base Refinement, Extraction des chronicles

\clearpage

\begin{center}
	\large \textbf{Abstract}
\end{center}

In the manufacturing domain, the detection of anomalies such as mechanical faults and failures enables the launching of predictive maintenance tasks, which aim to predict future faults, errors, and failures and also enable maintenance actions. With the trend of Industry 4.0, predictive maintenance tasks are benefiting from advanced technologies such as Cyber-Physical Systems (CPS), the Internet of Things (IoT), and Cloud Computing. These advanced technologies enable the collection and processing of sensor data that contain measurements of physical signals of machinery, such as temperature, voltage, and vibration.

However, due to the heterogeneous nature of industrial data, sometimes the knowledge extracted from industrial data is presented in a complex structure. Therefore formal knowledge representation methods are required to facilitate the understanding and exploitation of the knowledge. Furthermore, as the CPSs are becoming more and more knowledge-intensive, uniform knowledge representation of physical resources and reasoning capabilities for analytic tasks are needed to automate the decision-making processes in CPSs. These issues bring obstacles to machine operators to perform appropriate maintenance actions.

To address the aforementioned challenges, in this thesis, we propose a novel semantic approach to facilitate predictive maintenance tasks in manufacturing processes. In particular, we propose four main contributions: i) a three-layered ontological framework that is the core component of a knowledge-based predictive maintenance system; ii) a novel hybrid semantic approach to automate machinery failure prediction tasks, which is based on the combined use of chronicles (a more descriptive type of sequential patterns) and semantic technologies; iii) a new approach that uses clustering methods with Semantic Web Rule Language (SWRL) rules to assess failures according to their criticality levels; iv) a novel rule base refinement approach that uses rule quality measures as references to refine a rule base within a knowledge-based predictive maintenance system. These approaches have been validated on both real-world and synthetic data sets.

Keywords: Industry 4.0, Predictive Maintenance, Machinery Failure Prediction, Failure Criticality Assessment, Ontology Engineering, Rule-based Reasoning, Rule Base Refinement, Chronicle Mining


		%	\adjustmtc
			\dominitoc						% Génération des mini-toc	\pagenumbering{arabic}
			\tableofcontents
			
\chapter*{List of publications}
\addcontentsline{toc}{chapter}{List of publications}%
During this Ph.D, the following research papers have been published:
\begin{flushleft}
\textbf{Peer-reviewed journal papers}
\end{flushleft}

\begin{flushleft}
\textbf{Peer-reviewed international conference papers}
\end{flushleft}

		% Liste des figures
		%	\renewcommand*\listfigurename{Liste des figures}
			\listoffigures
		% Liste des tableaux
		\listoftables
%%%%%%%%%%%%%%%%%%%%%%%%%%%%%%%%%%%%%		
%        Contenu du document        %
%%%%%%%%%%%%%%%%%%%%%%%%%%%%%%%%%%%%%
	\setcounter{mtc}{2}	% "Corrige" les minitocs décallés à cause des chapter* (ex : table des matières)
	\pagenumbering{arabic}

	\part{Synthèse de la thèse en français}
	%%%%%%%%%%%%%%%%%%%%%%%%%%%%%%%%%%%%%%%%%%%%%%%%%%%%%%%%%%%%%%%%%%%%%%%%%%%%%%%%%%%%%%%%%%%%%
%%									Introduction											%
%%%%%%%%%%%%%%%%%%%%%%%%%%%%%%%%%%%%%%%%%%%%%%%%%%%%%%%%%%%%%%%%%%%%%%%%%%%%%%%%%%%%%%%%%%%%%
\begin{otherlanguage}{French}
\chapter{État de l'art}
\normalsize
\label{ch1fran}
\adjustmtc
\adjustmtc

\section{Contexte de la thèse}
\subsection{Contexte du projet}

\section{Motivation}
\subsection{Contributions de cette thèse}

\section{Maintenance prédictive dans le contexte de l'industrie 4.0}
\label{ch2fran}

\section{Modèles existants pour la maintenance prédictive industrielle}
\label{ch3fran}

\section{Systèmes fondés sur la connaissance : fondements théoriques}
\label{ch4fran}

\section{Systèmes de maintenance prédictive basés sur la connaissance: état de l'art}
\label{ch5fran}

\end{otherlanguage}



	%%%%%%%%%%%%%%%%%%%%%%%%%%%%%%%%%%%%%%%%%%%%%%%%%%%%%%%%%%%%%%%%%%%%%%%%%%%%%%%%%%%%%%%%%%%%%
%%									Chapitre 6												%
%%%%%%%%%%%%%%%%%%%%%%%%%%%%%%%%%%%%%%%%%%%%%%%%%%%%%%%%%%%%%%%%%%%%%%%%%%%%%%%%%%%%%%%%%%%%%
\begin{otherlanguage}{French}
\chapter{Contributions}
\normalsize
\label{ch6fran}

\section{Un cadre ontologique pour les systèmes de maintenance prédictive basés sur la connaissance}

\section{Une nouvelle approche sémantique hybride pour la maintenance prédictive}
\label{ch7fran}

\section{Évaluation des défaillances de machines sur la base des niveaux de criticité}
\label{ch8fran}

\section{Utiliser les mesures de qualité des règles pour l'élagage et l'intégration des règles}
\label{ch9fran}

\section{Le prototype du software: KSPMI}

\section{Conclusions et travaux futurs}
\subsection{Contributions de ce travail de thèse}
\subsection{Perspectives}

\end{otherlanguage}










	\part{Introduction and state of the art}
	%%%%%%%%%%%%%%%%%%%%%%%%%%%%%%%%%%%%%%%%%%%%%%%%%%%%%%%%%%%%%%%%%%%%%%%%%%%%%%%%%%%%%%%%%%%%%
%%									Introduction											%
%%%%%%%%%%%%%%%%%%%%%%%%%%%%%%%%%%%%%%%%%%%%%%%%%%%%%%%%%%%%%%%%%%%%%%%%%%%%%%%%%%%%%%%%%%%%%
\selectlanguage{english}
\setcounter{chapter}{2}
\chapter{Introduction}
\normalsize
\label{ch1}
	\minitoc
	\clearpage
	
This chapter gives a general introduction to this thesis. It starts with background information for this thesis, followed by the motivation for conducting this research work. After that, the main contributions of this thesis are presented. Finally, the thesis structure is presented in detail.

\section{Background}

\section{Motivation}

\section{Contributions of this thesis}

\section{Structure of the thesis}
	%%%%%%%%%%%%%%%%%%%%%%%%%%%%%%%%%%%%%%%%%%%%%%%%%%%%%%%%%%%%%%%%%%%%%%%%%%%%%%%%%%%%%%%%%%%%%
%%									Chapitre 1											%
%%%%%%%%%%%%%%%%%%%%%%%%%%%%%%%%%%%%%%%%%%%%%%%%%%%%%%%%%%%%%%%%%%%%%%%%%%%%%%%%%%%%%%%%%%%%%
\chapter{Predictive maintenance in the context of Industry 4.0}\label{ch2}
\normalsize
\minitoc
%%%%%%%%%%%%%%%%%%%%%%%%%%%%%%%%%%%%%%%%%%%%%%%%%%%%%%%%%%%%%%%%%%%%%%%%%%%%%%%%%%%%%%%%%%%%%
\clearpage

	%%%%%%%%%%%%%%%%%%%%%%%%%%%%%%%%%%%%%%%%%%%%%%%%%%%%%%%%%%%%%%%%%%%%%%%%%%%%%%%%%%%%%%%%%%%%%
%%									Chapitre 3											%
%%%%%%%%%%%%%%%%%%%%%%%%%%%%%%%%%%%%%%%%%%%%%%%%%%%%%%%%%%%%%%%%%%%%%%%%%%%%%%%%%%%%%%%%%%%%%

\chapter{Existing models for industrial predictive maintenance}
\normalsize
\label{ch3}

\minitoc
\clearpage
	%%%%%%%%%%%%%%%%%%%%%%%%%%%%%%%%%%%%%%%%%%%%%%%%%%%%%%%%%%%%%%%%%%%%%%%%%%%%%%%%%%%%%%%%%%%%%
%%									Chapitre 4											%
%%%%%%%%%%%%%%%%%%%%%%%%%%%%%%%%%%%%%%%%%%%%%%%%%%%%%%%%%%%%%%%%%%%%%%%%%%%%%%%%%%%%%%%%%%%%%

\chapter{Knowledge-based systems: theoretical foundations}\label{ch4}
\normalsize
\minitoc

\clearpage
	%%%%%%%%%%%%%%%%%%%%%%%%%%%%%%%%%%%%%%%%%%%%%%%%%%%%%%%%%%%%%%%%%%%%%%%%%%%%%%%%%%%%%%%%%%%%%
%%									Chapitre 5											%
%%%%%%%%%%%%%%%%%%%%%%%%%%%%%%%%%%%%%%%%%%%%%%%%%%%%%%%%%%%%%%%%%%%%%%%%%%%%%%%%%%%%%%%%%%%%%
\chapter{Knowledge-based predictive maintenance systems: state of the art}\label{ch5}
\normalsize
	\minitoc
	
	\clearpage
	\part{Contributions}
	%%%%%%%%%%%%%%%%%%%%%%%%%%%%%%%%%%%%%%%%%%%%%%%%%%%%%%%%%%%%%%%%%%%%%%%%%%%%%%%%%%%%%%%%%%%%%
%%									Chapitre 6												%
%%%%%%%%%%%%%%%%%%%%%%%%%%%%%%%%%%%%%%%%%%%%%%%%%%%%%%%%%%%%%%%%%%%%%%%%%%%%%%%%%%%%%%%%%%%%%

\chapter{An ontological framework for knowledge-based predictive maintenance systems}
\normalsize
\label{ch6}
	\minitoc

%%%%%%%%%%%%%%%%%%%%%%%%%%%%%%%%%%%%%%%%%%%%%%%%%%%%%%%%%%%%%%%%%%%%%%%%%%%%%%%%%%%%%%%%%%%%%
\clearpage
	%%%%%%%%%%%%%%%%%%%%%%%%%%%%%%%%%%%%%%%%%%%%%%%%%%%%%%%%%%%%%%%%%%%%%%%%%%%%%%%%%%%%%%%%%%%%%
%%									Chapitre 7												%
%%%%%%%%%%%%%%%%%%%%%%%%%%%%%%%%%%%%%%%%%%%%%%%%%%%%%%%%%%%%%%%%%%%%%%%%%%%%%%%%%%%%%%%%%%%%%

\chapter{A novel hybrid semantic approach for predictive maintenance}
\normalsize
\label{ch7}
	\minitoc
\clearpage
	%%%%%%%%%%%%%%%%%%%%%%%%%%%%%%%%%%%%%%%%%%%%%%%%%%%%%%%%%%%%%%%%%%%%%%%%%%%%%%%%%%%%%%%%%%%%%
%%									Chapitre 8												%
%%%%%%%%%%%%%%%%%%%%%%%%%%%%%%%%%%%%%%%%%%%%%%%%%%%%%%%%%%%%%%%%%%%%%%%%%%%%%%%%%%%%%%%%%%%%%

\chapter{Assessment of machinery failures based on criticality levels}
\normalsize
\label{ch8}
\minitoc
\clearpage
	%%%%%%%%%%%%%%%%%%%%%%%%%%%%%%%%%%%%%%%%%%%%%%%%%%%%%%%%%%%%%%%%%%%%%%%%%%%%%%%%%%%%%%%%%%%%%
%%									Chapitre 9												%
%%%%%%%%%%%%%%%%%%%%%%%%%%%%%%%%%%%%%%%%%%%%%%%%%%%%%%%%%%%%%%%%%%%%%%%%%%%%%%%%%%%%%%%%%%%%%

\chapter{Using rule quality measures for rule base pruning and integration}
\normalsize
\label{ch9}
	\minitoc
\clearpage
	%%%%%%%%%%%%%%%%%%%%%%%%%%%%%%%%%%%%%%%%%%%%%%%%%%%%%%%%%%%%%%%%%%%%%%%%%%%%%%%%%%%%%%%%%%%%%
%%									Chapitre 10												%
%%%%%%%%%%%%%%%%%%%%%%%%%%%%%%%%%%%%%%%%%%%%%%%%%%%%%%%%%%%%%%%%%%%%%%%%%%%%%%%%%%%%%%%%%%%%%

\chapter{The software prototype: KSPMI}
\normalsize
\label{ch10}
	\minitoc
\clearpage

	%%%%%%%%%%%%%%%%%%%%%%%%%%%%%%%%%%%%%%%%%%%%%%%%%%%%%%%%%%%%%%%%%%%%%%%%%%%%%%%%%%%%%%%%%%%%%
%%									Conclusion											%
%%%%%%%%%%%%%%%%%%%%%%%%%%%%%%%%%%%%%%%%%%%%%%%%%%%%%%%%%%%%%%%%%%%%%%%%%%%%%%%%%%%%%%%%%%%%%
\chapter{Conclusions and future work}
\normalsize
\minitoc
%%%%%%%%%%%%%%%%%%%%%%%%%%%%%%%%%%%%%%%%%%%%%%%%%%%%%%%%%%%%%%%%%%%%%%%%%%%%%%%%%%%%%%%%%%%%%
\clearpage
\label{Conclusions and future work}


	
	
	% Annexes
	\part{Appendix}
	\begin{appendix}
	%	\pagenumbering{Roman}
                %%%%%%%%%%%%%%%%%%%%%%%%%%%%%%%%%%%%%%%%%%%%%%%%%%%%%%%%%%%%%%%%%%%%%%%%%%%%%%%%%%%%%%%%%%%%%
%%									ANNEXES 												%
%%%%%%%%%%%%%%%%%%%%%%%%%%%%%%%%%%%%%%%%%%%%%%%%%%%%%%%%%%%%%%%%%%%%%%%%%%%%%%%%%%%%%%%%%%%%%
\chapter{The fuzzy \textit{c}-means clustering algorithm}
\normalsize
\label{The fuzzy c-means clustering algorithm}

                %%%%%%%%%%%%%%%%%%%%%%%%%%%%%%%%%%%%%%%%%%%%%%%%%%%%%%%%%%%%%%%%%%%%%%%%%%%%%%%%%%%%%%%%%%%%%
%%									ANNEXES 												%
%%%%%%%%%%%%%%%%%%%%%%%%%%%%%%%%%%%%%%%%%%%%%%%%%%%%%%%%%%%%%%%%%%%%%%%%%%%%%%%%%%%%%%%%%%%%%
\chapter{The theory of evidential clustering}
\normalsize
\label{The evidential c-means clustering algorithm}
%%%%%%%%%%%%%%%%%%%%%%%%%%%%%%%%%%%%%%%%%%%%%%%%%%%%%%%%%%%%%%%%%%%%%%%%%%%%%%%%%%%%%%%%%%%%%
This Appendix presents the theory of evidential \textit{c}-means (ECM) clustering algorithm. The following basic notions and theory are taken from the research paper \cite{masson2008ecm}.


                %%%%%%%%%%%%%%%%%%%%%%%%%%%%%%%%%%%%%%%%%%%%%%%%%%%%%%%%%%%%%%%%%%%%%%%%%%%%%%%%%%%%%%%%%%%%%
%%									ANNEXES 												%
%%%%%%%%%%%%%%%%%%%%%%%%%%%%%%%%%%%%%%%%%%%%%%%%%%%%%%%%%%%%%%%%%%%%%%%%%%%%%%%%%%%%%%%%%%%%%
\chapter{The fast non-dominated sorting algorithm}
\normalsize
\label{The fast non-dominated sorting procedure}
%%%%%%%%%%%%%%%%%%%%%%%%%%%%%%%%%%%%%%%%%%%%%%%%%%%%%%%%%%%%%%%%%%%%%%%%%%%%%%%%%%%%%%%%%%%%%
This Appendix gives the theory of multi-objective optimization and the fast non-dominated sorting algorithm. The theory is taken from research papers \cite{ortega2017non} and \cite{miettinen2012nonlinear}.

	\end{appendix}
  %    \pagenumbering{Roman}
      \bibliographystyle{alpha}
      \bibliography{These/Chapitres/Manuscrit.bib}
      %\printbibliography[title={Reference}]
      
	
\end{document}
